\chapter{Sicherheitskonzept}
\label{ch:sicherheit}

% ==============================================================================
\section{Gefährdungen durch eine gelernte Policy}
\label{sec:gefahren}

Eine programmierte Bahn lässt sich vor der Fahrt vollständig prüfen. Bei einer gelernten
Policy ist das nicht möglich: Sie erzeugt ihre Sollwerte erst während der Fahrt, und in
Zuständen außerhalb der Trainingsdaten ist ihr Verhalten nicht vorhersagbar
(\cref{sec:covariate-shift}). Typische Fehlbilder aus den Versuchen sind ein Abtauchen in
Richtung Tisch, Sprünge zwischen zwei Vorhersagen und ein Stehenbleiben an Haltepunkten.
Hinzu kommen technische Fehler wie ein eingefrorenes Kamerabild oder ein unterbrochener
Steuerkanal.

NeuraPy bietet keine frei definierbaren Schutzräume. Modul~4 baut deshalb eine eigene,
mehrstufige Sicherheitsschicht zusätzlich zu den Schutzfunktionen des Roboters
(\cref{fig:sicherheit}). Der Hardware-Not-Halt bleibt bei allen Fahrten in Reichweite; ein
Software-Stopp ersetzt ihn nicht.

\begin{figure}[htbp]
  \centering
  \begin{tikzpicture}[font=\small, node distance=3.5mm,
      st/.style={draw, rounded corners=1.5mm, align=left, text width=96mm, minimum height=9mm,
                 inner xsep=3mm},
      tag/.style={font=\scriptsize\bfseries, anchor=east, align=right, text width=26mm},
      arr/.style={-{Stealth[length=2mm]}, thick}]
    \node[st, fill=blue!6] (s1) {Weiche Begrenzung: Änderung je Takt $\le$ \SI{0.8}{rad/s}, Richtung bleibt};
    \node[st, fill=blue!10, below=of s1] (s2) {Zielprüfung: Ziel und Nachlauf gegen Tisch, Sperrzonen,
      Hindernisse und Hülle der Demonstrationen};
    \node[st, fill=blue!14, below=of s2] (s3) {Harte Grenze je \SI{60}{\hertz}-Sollwert: endlich, Achsgrenzen,
      $\le$ \SI{1.0}{rad/s}, Ziel--Ist $\le$ \SI{0.35}{rad}};
    \node[st, fill=orange!12, below=of s3] (s4) {Überwachung (eigener Thread, \SI{20}{\hertz}): gemessene
      Stellung gegen Kollisionsmodell};
    \node[st, fill=orange!16, below=of s4] (s5) {Bildalter, stehender Arm, Vorhersage zu spät};
    \node[st, fill=red!10, below=of s5] (s6) {Schutzfunktionen der Steuerung, Hardware-Not-Halt};
    \node[tag] at ([xshift=-2mm]s1.west) {vor dem Senden};
    \node[tag] at ([xshift=-2mm]s4.west) {während der Fahrt};
    \node[tag] at ([xshift=-2mm]s6.west) {Anlage};
    \draw[arr] ([xshift=4mm]s1.east) -- ([xshift=4mm]s6.east)
      node[midway, right, font=\scriptsize, align=left] {letzte\\Instanz};
  \end{tikzpicture}
  \caption[Stufen der Sicherheitsschicht]{Stufen der Sicherheitsschicht von Modul~4. Die
  oberen Stufen prüfen jeden Befehl, bevor er die Steuerung erreicht, die unteren den
  tatsächlichen Zustand während der Fahrt.}
  \label{fig:sicherheit}
\end{figure}

% ==============================================================================
\section{Prüfung vor dem Senden}
\label{sec:pruefung-senden}

\paragraph{Weiche Begrenzung.} Die Änderung der Sollstellung je Takt wird auf eine
Gelenkgeschwindigkeit von \SI{0.8}{rad/s} begrenzt; die Richtung bleibt erhalten.
\emph{Begründung:} Die Daten enthalten Korrekturen bis etwa \SI{1}{rad/s}, und beim
Eintreffen einer verspäteten Vorhersage springt der Mittelwert kurzzeitig. Ohne diese Stufe
stoppte die harte Grenze eine sonst gute Fahrt bei \SI{1.03}{rad/s}.

\paragraph{Zielprüfung.} Jedes Ziel der Policy wird vor dem Senden gegen das Kollisionsmodell
geprüft. Dazu werden aus den Gelenkwinkeln die Positionen von Werkzeug, Kamera, Handgelenk und
Ellbogen berechnet und mit Tisch, Sperrzonen und festen Hindernissen verglichen. Der Tisch
wird live aus den Tischpunkten der Steuerung gelesen, ersatzweise aus den Werten der Aufnahme.
Zusätzlich muss das Werkzeug innerhalb der um \SI{3}{\centi\metre} erweiterten Hülle aller
Demonstrationen bleiben.

\begin{description}
  \item[Problem] Zuerst prüfte nur die Überwachung die \emph{gemessene} Stellung. Sie reagiert
    erst, wenn der Arm schon an der Grenze ist; mit Prüftakt und Stoppzeit kommt der Stopp
    einige Millimeter zu spät. Und nach einem abgelehnten Ziel lief der Arm an der Anlage noch
    bis zu \SI{7.5}{\milli\metre} nach.
  \item[Lösung] Geprüft wird das Ziel selbst, und zusätzlich das um einen Takt linear
    verlängerte Ziel. So wird bereits abgelehnt, solange der Nachlauf bei der aktuellen
    Geschwindigkeit noch vor der Grenze endet.
\end{description}

\paragraph{Harte Grenze.} Jeder einzelne \SI{60}{\hertz}-Sollwert wird im Adapter verworfen,
wenn er nicht endlich ist, außerhalb der Achsgrenzen liegt oder eine Gelenkgeschwindigkeit
über \SI{1.0}{rad/s} gegenüber dem letzten Sollwert erfordert. Die Zeitbasis dafür ist nach
unten begrenzt, ein Stocken des Takts erlaubt also keinen großen Sprung. Liegt die
Sollstellung mehr als \SI{0.35}{rad} von der gemessenen entfernt, wird ebenfalls nichts
gesendet. In allen Fällen stoppt der Roboter. Diese Grenze gilt auch bei Aufnahme und Teachen.

% ==============================================================================
\section{Überwachung während der Fahrt}
\label{sec:ueberwachung}

\begin{description}
  \item[Arbeitsraum] Ein eigener Thread prüft alle \SI{50}{\milli\second} die gemessene
    Stellung gegen dasselbe Kollisionsmodell. Er läuft getrennt von der Vorhersage, damit eine
    verzögerte Vorhersage den Stopp nicht verzögert.
  \item[Eingefrorenes Bild] Ist ein Kamerabild zwei Takte in Folge älter als
    \SI{100}{\milli\second}, hält die Fahrt an. Die Policy würde sonst auf ein veraltetes Bild
    reagieren, ohne es zu merken.
  \item[Stehender Arm] Bewegt sich der Arm fünf Takte lang praktisch nicht, obwohl das Ziel
    mehr als \SI{0.02}{rad} entfernt ist, endet die Fahrt. \emph{Anlass:} Der Not-Halt am
    Bedienhandgerät hält den Arm an, ohne dass \code{servo_j} einen Fehler meldet. Die Fahrt
    schickte weiter Ziele, und nach dem Entriegeln hätte ein weit entferntes Ziel angelegen.
  \item[Vorhersage zu spät] Liegt keine gültige Vorhersage mehr für den aktuellen Takt vor,
    endet die Fahrt (\cref{sec:takt}).
\end{description}

% ==============================================================================
\section{Stoppverhalten und Freigabe}
\label{sec:stopp}

Der Software-Stopp hielt den Arm an der Anlage innerhalb von \SI{96}{\milli\second} an. Der
Befehl zum Beenden der Online-Bewegung bremst dagegen nicht, er wird deshalb für Abbrüche
nicht verwendet. Fahrten am realen Roboter erfordern eine ausdrückliche Bestätigung; aus dem
Ablauf des Gesamtsystems wird die Freigabe des dortigen Sicherheitsdialogs übernommen
(\cref{ch:integration}).

\paragraph{Grenzen.} Die Sicherheitsschicht erkennt, was sich im Kollisionsmodell abbilden
lässt. Ein Greifen neben das Bauteil oder ein Anstoßen des Bauteils innerhalb des erlaubten
Raums verhindert sie nicht. Die Grenzwerte der Gelenkgeschwindigkeiten stammen aus Daten der
\ac{vm} und wurden an der Anlage bestätigt; die Achsgrenzen stammen aus dem \ac{urdf}.

\todo[inline]{Achsgrenzen gegen Datenblatt bzw.\ Bedienhandgerät abgleichen (im Code noch als
nicht verifiziert markiert).}
